\section{Introduction}
\label{sec:intro}

A selective classifier returns a decision only when its confidence clears a
cutoff and refers the remaining cases to a human reader. The cutoff fixes an
operating point, a coverage, and with it the workload the system places on
radiology. In practice the cutoff, like the per-finding decision thresholds, is
chosen on data from the hospital where the system was developed. If the
system is then deployed elsewhere unchanged, the operational question is
whether that frozen policy still delivers the error rate and the referral rate
it was calibrated for. Re-fitting the cutoff at the new site answers a related
but different question: whether the method can be re-tuned.

Multimodal chest-radiograph models add a second source of shift. Models that
read the clinical indication or history alongside the image depend on text
whose availability and content reflect local documentation practice. A model
that relies on that text at one site may meet thinner, absent or differently
written context at the next. Institutional and contextual shift therefore
travel together, and an evaluation that varies only one of them cannot
estimate how they combine.

We designed and preregistered an evaluation that crosses the two. A source
hospital (MIMIC-CXR-JPG) was used for development and for freezing the policy,
and an external hospital (CheXpert Plus) for evaluation only. At each site
pre-diagnostic context was left intact, removed, or replaced with another
patient's. The policy was transferred without refitting.

This revision reports that evaluation together with an audit carried out after
the external results had been inspected, and the audit changed the principal
conclusion. The selective-risk estimator used in the original analysis counted
studies whose report mentioned none of the five target pathologies as
correctly classified. Such studies were far more common at the source than at
the external site, and this produced the reported cross-site degradation.
Correcting the defect removes that degradation under the primary label
handling. The corrected cross-site
contrast is not stable either: it moves by more than its own interval across
defensible choices about uncertain and unmentioned labels, aggregation, label
source, training seed and view stratum. We therefore report the study as two
linked contributions:
\begin{itemize}
  \item a \emph{compound-shift evaluation protocol}, specified as a reusable
        machine-readable registry with frozen thresholds, controlled context
        interventions and explicit invariants (\cref{sec:methods}), together
        with its worked application; and
  \item a \emph{controlled demonstration that the cross-site verdict is
        endpoint-dependent}. Holding studies, predictions, thresholds and
        cutoffs fixed and changing one factor at a time, we show which
        analysis choices drive the sign of the transfer estimate and why
        (\cref{sec:res-label-source}).
\end{itemize}
Within-site findings about context interventions are reported alongside. For
the late-fusion model at the source site they held across every specification
and both training seeds. At the external site they did not survive the second
seed.

We claim no new architecture or uncertainty method; the backbones are
conventional and held fixed. This is not a prospective, clinical or regulatory
validation, and the labels are automatic report-derived labels, not an expert
reference standard (\cref{sec:limitations}). Every analysis added after the
external results were seen is marked as post hoc, dated, and reported
alongside the original estimates, which are retained.
